\section{Conclusion and Future Work}
\label{sec:conclusion}

In this paper, we presented the first systematic investigation of \textit{over-privileged tool selection} under cross-lingual transfer into regional Southeast Asian languages (Indonesian and Javanese). We introduced \textsc{ToolPrivBench-SEA}, showing that morphological and honorific ambiguity in regional dialects exacerbates the tendency of LLM agents to invoke excessively permissive tools.

To resolve this vulnerability, we designed and evaluated a deterministic, lattice-based least-privilege verifier. Our guardrail achieves three critical properties:
\begin{enumerate}
    \item \textbf{Provable Least Privilege:} Eliminates over-privileged API calls ($OPR \rightarrow 0.0\%$).
    \item \textbf{Transient Fault Resilience:} Prevents panic privilege escalation when standard tools encounter timeouts.
    \item \textbf{Negligible Overhead:} Operates with sub-15ms latency without requiring model fine-tuning.
\end{enumerate}

\paragraph{Future Work}
We plan to expand our framework into decentralized multi-agent environments, exploring how cooperative multi-agent workflows handle cascading permissions, and scaling the benchmark across additional Austronesian languages, including Sundanese and Tagalog.
